\documentclass[12pt]{article}
\input{.house-style/preamble.tex}
\usepackage{xurl}
\usepackage{float}
\setlength{\headheight}{14pt}
\clubpenalty=10000
\widowpenalty=10000
\emergencystretch=1em
\AtBeginBibliography{\emergencystretch=1em}
\setcounter{biburlnumpenalty}{100}
\newcommand{\tablebody}[1]{\csname @@input\endcsname #1 }
\hypersetup{pdftitle={Coverage audit: how the four implementations cover the fragment's constructions}}
\title{Coverage audit: how the four implementations cover the fragment's constructions}
\author{Brett Reynolds}
\date{Supplementary material, September 2026}
\fancyhead[L]{\small\scshape Coverage audit}
\usepackage{longtable}
\begin{document}
\maketitle

\section{What this audit is}

The article's §5 compares four implementations of the same facts: the D-noun categorization with ordinary Head (I1), a separate determinative category with ordinary Head (I2), a separate category with \textit{CGEL}'s fusion of functions (I3), and the D-noun categorization with fusion (I4). It states which conditions the accounts share and which are specific to one of them. This audit records, for each construction family the fragment covers, and for each lexeme group the article distinguishes, it records how each implementation covers the facts: by a general rule of that grammar (D), by a general rule that applies only because the grammar admits a determinative as the lexical head of Nom and NP (Dn), by a rule the grammar has to state once for nouns and again for determinatives, or once for NP and again for DP (D2), by a fact stored on the lexeme (L), by a condition stated for that construction alone (S), or not at all (U).

The classifications are the author's, made by hand from the fragment's stated conditions and the sections cited beside each cell; they are the analysis, not data. The precedence rule is this: a cell records the mechanism by which the implementation covers the licensed use. Lexical permissions, target selection, forms, and inflection are held fixed across the four implementations and are coded L only where the fact is nothing but such a permission or exclusion; where a licensed use is covered by a rule, the rule's class is recorded even though the permission is presupposed. S and D2 codes count cells; the distinct conditions and duplicated rules behind them are listed separately (Table~\ref{tab:rules}), since a condition stated once may apply in several cells. The accounts are compared as complete statements of the same fragment, over the rules and facts each has to state, and not by the number of categories each posits; that is the standard the minimum description length principle makes explicit \citep{rissanen1978mdl}. Beside each cell the audit prints the number of participation claims the claim register records for that family and group, as licensed/conditional/excluded, so that a cell's classification can be checked against what the article asserts there. Cells with no claim and no analysis in the article are left out; the article's analysis, not the register, decides which cells exist.

\section{The four implementations and their primitives}

The classes count rules, not primitives. The primitives differ too, and Table~\ref{tab:prim} lists them, since a rule that is general in one grammar may be general only because that grammar has a primitive the other lacks.

\begin{table}[H]\centering\small
\caption{Primitives of the four implementations.}\label{tab:prim}
\begin{tabular}{>{\raggedright\arraybackslash}p{4.4cm}>{\raggedright\arraybackslash}p{2.3cm}>{\raggedright\arraybackslash}p{2.3cm}>{\raggedright\arraybackslash}p{2.3cm}>{\raggedright\arraybackslash}p{2.3cm}}
\toprule
 & I1 & I2 & I3 & I4 \\
\midrule
Lexical categories & Noun, four subcategories & Noun (three) and D & Noun (three) and D & Noun, four subcategories \\
Category condition on lexical Heads in Nom & N & N or D & N or D & N \\
Phrase types in Det & NP, PP & NP, PP & DP, NP, PP & NP, PP \\
Head configurations for independent determinatives & ordinary Head & ordinary Head & Det--Head and Mod--Head fusion & Det--Head and Mod--Head fusion \\
Adjectival Mod--Head fusion (\mention{the rich}) & yes & yes & yes & yes \\
Restrictions stated on the fused inner phrase & none & none & yes & yes \\
Complex determinatives listed & none & none & \mention{a few}, \mention{a little}, \mention{many a}, \mention{a great many} & none \\
Conversion rule for the set-denoting use & none & yes & yes & none \\
Domain of the head-genitive definition & Noun heads & Noun and D heads & Noun and D heads & Noun heads \\
\bottomrule
\end{tabular}
\end{table}

\section{Results}

\tablebody{analysis/generated/coverage-summary.tex}

Fusion is held fixed for adjectives (\mention{the rich}) in all four implementations. \textit{CGEL} (I3) uses it a second time, for independent determinatives, and pays for that use with one condition on the fused inner phrase (nine cells) and with rules stated once for NPs and again for DPs (twelve cells, four rules). The two ordinary-Head implementations remove the second use and share the saving; they buy it differently. D-noun (I1) moves the determinatives one node down the taxonomy and keeps the head condition as it is. Separate D with ordinary Head (I2) keeps D a primary category and adds a disjunct to the head condition, and the Dn column records the 32 derivations that go through that disjunct. For the same coverage, I2 therefore carries what I1 carries plus the disjunct, plus a primary category whose defining profile duplicates Noun's (the nominal connections of the article's §2), plus the conversion rule and the extended genitive definition (Table~\ref{tab:rules}). A cross-cutting nominal feature wouldn't change the count. A feature carried by exactly the nouns and the determinatives has the extension of the superordinate node, and it costs its definition and one assignment per lexeme where the node costs itself and its inheritance links; a feature would earn its place by cross-cutting the categories, and in this fragment none does. The lexical facts (use permissions, target selection, forms, inflection) are the same in all four.

The sensitivity tables apply the second-round report's alternative principles: lexical licensing first, the inner-phrase restriction counted once, typed schemas in place of duplicated rules, the secondary use as a rule in every account. Under those principles I1, I3, and I4 come out level on cells and I2 trails only by its Dn cells; pricing the disjunct at zero brings I2 level too. That is where the disagreement lies: the recoding prices the disjunct at zero and doesn't count the category it keeps. The distinct-conditions table counts both, and there I2 is I1 plus one condition and one category. The article's §5.5 states the comparison in these terms.

\section{The grid}

\tablebody{analysis/generated/coverage-audit.tex}

\section{Limits}

The grid covers the fragment's constructions and nothing outside them: predication, interrogative clause structure, coordination, and modification outside NP structure are not classified, as the article's §5.1 says of the fragment itself. The classes are coarse by design; a D under one implementation and a D under another can rest on different primitives (Table~\ref{tab:prim}), and the summary counts should be read with that table. The classifications are contestable at the cell, and the JSON file (\texttt{analysis/coverage-audit.json}) is written so that a reader can change one cell and regenerate the tables (\texttt{analysis/tools/coverage\_audit.py}).

\printbibliography

\end{document}
